\section{De dyraste misstagen}
\label{sec:mistakes}

Var och en av dem har kostat någon en kväll. De står i den ordning de brukar inträffa.

\begin{booktable}{@{}rLL@{}}
\thead{\#} & \thead{Misstag} & \thead{Symptom} \\ \midrule
1 & Parsern kontrollerar inte längdfältet mot buffertens storlek & Buffertöverskrivning; allt
  möjligt annat går sönder \\
2 & Parsern går tillbaka till \code{WaitForSof1} på ett andra \code{0xA5} & Enstaka frames
  missas, oförklarligt \\
3 & Flerbytesfält serialiseras genom att kopiera minnet & Fungerar tills den andra sidan byggs
  med annan endianness \\
4 & Dubblett kvitteras inte igen & Sändaren skickar om i all evighet \\
5 & \code{TX_DATA_LO} och \code{TX_DATA_HI} förväxlade & Framen kommer fram --- bakvänd \\
6 & Cast efter skiftningen i \code{readReg()} & Fungerar för små värden, går sönder vid bit 7 \\
7 & \code{TX_SEND} skrivs före dataregistren & Noden ligger exakt en frame efter \\
8 & \code{clearError()} skriver \code{0x1} & Felbiten går aldrig ned; all felhantering slutar
  fungera \\
9 & \code{RX_ACK} glöms & Samma frame läses i en oändlig loop \\
10 & \code{receive()} maskerar inte mot \code{dlc} & Upp till sex bytes rapporteras som aldrig
  sändes \\
11 & \code{while (!send(frame)) {}} på en ogiltig frame & Hänger för alltid \\
12 & Pinnmuxen inte satt & Programmet kör, ledningarna är tysta, inga fel rapporteras \\
13 & Dataregistret läses inte i \code{transfer()} & All data ligger en byte fel \\
14 & \code{VDDIO2} inte kopplad & Ser ut precis som en trasig SPI-implementation \\
15 & Factoryns medlemmar i fel deklarationsordning & Referens till ett objekt som inte
  konstruerats \\
\end{booktable}

\begin{aside}
\note Elva av de femton hade fångats av ett test på laptopen. Tre av dem --- 12, 13 och 14 ---
hade det inte, och det är precis de tre som bor i den enda klass inget automattest täcker. Det är
ingen slump, och det är hela argumentet för stegen i \kapref{ch:bringup}.
\end{aside}

\section{Efter kursen}

Två saker fortsätter direkt härifrån.

\textbf{Testning.} Genom hela projektet har du kört och läst tester som någon annan skrivit. Nästa
steg är att bestämma vad de ska kontrollera --- vilket är ett svårare problem än det låter, och
det som avgör om en testsvit är värd något.

\textbf{Ditt eget repo.} Drivrutinen du byggt är en komplett, testad, lagerdelad drivrutinsstack
mot riktig hårdvara, utvecklad mot ett kontrakt och verifierad utan tillgång till den andra halvan.
Det är ett rimligt första objekt i en portfölj, och det är en mer realistisk beskrivning av
inbyggd utveckling än de flesta kursprojekt.
